\documentclass[11pt,a4paper]{article}
\usepackage[margin=1in]{geometry}
\usepackage{amsmath,amsthm,amssymb}
\usepackage{mathtools}
\usepackage{array}
\usepackage[hidelinks]{hyperref}

\newtheorem{theorem}{Theorem}
\renewcommand{\thetheorem}{\Alph{theorem}}
\newtheorem{lemma}{Lemma}
\newtheorem{proposition}{Proposition}
\newtheorem{corollary}{Corollary}
\theoremstyle{definition}
\newtheorem{definition}{Definition}
\theoremstyle{remark}
\newtheorem{remark}{Remark}

\newcommand{\R}{\mathbb{R}}
\newcommand{\T}{\mathbb{T}}
\newcommand{\Z}{\mathbb{Z}}
\newcommand{\nablau}{\nabla u}
\newcommand{\abs}[1]{\lvert#1\rvert}
\newcommand{\norm}[1]{\lVert#1\rVert}
\newcommand{\dd}{\mathrm{d}}
\newcommand{\pv}{\mathrm{P.V.}}

\hypersetup{
  pdftitle={Direction of vorticity and continuation for 3D Navier-Stokes: the 1/2-Holder criterion on the maximal classical interval},
  pdfauthor={Benjamin Stanley Frohman}
}

\title{Direction of vorticity and continuation for 3D Navier--Stokes:\\
the \(\tfrac12\)-H\"older criterion on the maximal classical interval}
\author{Benjamin Stanley Frohman%
\thanks{No tether term is used. Degeneracy of
\(\xi=\omega/\abs{\omega}\) is respected. Theorem~\ref{thm:A} is the
Beir\~ao da Veiga--Berselli implication. Theorem~\ref{thm:B} records
the result of constructing \(\Omega,C\) from the unmodified PDE: the
a priori bound on \(C\) does not close, so neither global smoothness
nor blow-up is obtained for arbitrary data.}}
\date{1 September 2026}

\begin{document}
\maketitle

\begin{abstract}
Let \(u\) be the unique classical solution of the 3D incompressible
Navier--Stokes equations on the maximal interval \([0,T^*)\), viscosity
\(\nu>0\). Write \(\omega=\nabla\times u\), \(r=\abs{\omega}\),
\(\xi=\omega/r\) on \(\{r>0\}\), and
\(S=\tfrac12\bigl(\nabla u+(\nabla u)^\top\bigr)\).
The two constants of the criterion are
\[
\Omega=\Omega(u_0,\nu)>0,\qquad C=C(u_0,\nu)<\infty.
\]
If
\[
\abs{\xi(x,t)-\xi(y,t)}\le C\,\abs{x-y}^{1/2}
\]
whenever \(r(x,t),r(y,t)\ge\Omega\) and \(t<T^*\), then \(T^*=\infty\)
and \(u\) is globally smooth. The direction \(\xi\) is never evaluated
on \(\{r=0\}\). No additional geometric term is added to the PDE.
The implication is the Beir\~ao da Veiga--Berselli \(\tfrac12\)-H\"older
form of the Constantin--Fefferman criterion. An attempt to construct
\(\Omega\) and \(C\) from~\eqref{eq:NS}--\eqref{eq:div} for arbitrary
smooth divergence-free \(u_0\) is carried out below. The resulting
a priori bound on \(C\) does not close at the energy level. The test
therefore does not decide global smoothness versus blow-up.
\end{abstract}

\noindent\textit{MSC-style keywords:} Navier--Stokes, vorticity direction,
Beale--Kato--Majda, Constantin--Fefferman, Beir\~ao da Veiga--Berselli.

\tableofcontents

\section{Scope}

This note treats the unmodified Cauchy problem for 3D incompressible
Navier--Stokes, written as~\eqref{eq:NS}--\eqref{eq:div} below, on the
maximal classical existence interval \([0,T^*)\). The only geometric
input is a uniform spatial \(\tfrac12\)-H\"older bound on the vorticity
\emph{direction} on the set where the vorticity \emph{magnitude} is at
least a positive threshold. From that input one concludes \(T^*=\infty\).

The following are \textbf{not} proved: that \(C(t)\) remains finite up to
\(T^*\) for every smooth divergence-free \(u_0\); that \(T^*<\infty\)
occurs for some such \(u_0\); that a named tether or metriplectic
correction is present in the equations. The PDE used is~\eqref{eq:NS}.

Section~\ref{sec:test} constructs candidate constants from the PDE and
tests whether the energy identity forces them to stay finite. It does not.

\section{The equations}

Let \(\nu>0\). On \(\T^3=\R^3/(2\pi\Z)^3\) (the whole-space case is
identical up to replacing the periodized Newtonian potential by
\(1/(4\pi\abs{x})\)) one considers
\begin{align}
\partial_t u+(u\cdot\nabla)u&=\nu\Delta u-\nabla p,
\label{eq:NS}\\
\nabla\cdot u&=0,
\label{eq:div}
\end{align}
with \(u(\,\cdot\,,0)=u_0\in C^\infty(\T^3;\R^3)\), \(\nabla\cdot u_0=0\).
Local well-posedness in \(H^s\), \(s>5/2\), yields a unique classical
solution on a maximal interval \([0,T^*)\), \(T^*\in(0,\infty]\), with
the standard blow-up alternative: if \(T^*<\infty\) then
\begin{equation}
\label{eq:H1blowup}
\limsup_{t\uparrow T^*}\norm{\nabla u(t)}_{L^2(\T^3)}=\infty.
\end{equation}
Equivalently, \(\limsup_{t\uparrow T^*}\norm{\omega(t)}_{L^2}=\infty\),
since \(\norm{\omega}_{L^2}=\norm{\nabla u}_{L^2}\) under~\eqref{eq:div}.
This is the energy-class continuation criterion (Kato); it is the only
continuation tool used below. (Beale--Kato--Majda in \(L^1_tL^\infty_x\)
is stronger than needed once enstrophy is bounded.)

Curl of~\eqref{eq:NS}, using~\eqref{eq:div}, gives the vorticity equation
\begin{equation}
\label{eq:vort}
\partial_t\omega+(u\cdot\nabla)\omega=(\omega\cdot\nabla)u+\nu\Delta\omega,
\qquad\omega=\nabla\times u.
\end{equation}
Write \(r:=\abs{\omega}\) and, on the open set \(\{r>0\}\),
\begin{equation}
\label{eq:xi}
\xi:=\frac{\omega}{r}.
\end{equation}
The strain is the symmetric gradient
\begin{equation}
\label{eq:S}
S:=\frac12\bigl(\nabla u+(\nabla u)^\top\bigr).
\end{equation}
On \(\{r>0\}\) one has the stretching identity
\begin{equation}
\label{eq:stretch}
(\omega\cdot\nabla)u\cdot\omega=r^2\,(S\xi\cdot\xi).
\end{equation}
The left-hand side of~\eqref{eq:stretch} extends continuously by \(0\)
at \(r=0\). The unit field \(\xi\) is \emph{not} extended; every estimate
below that mentions \(\xi(x)\) is restricted to \(r(x)>0\).

\section{The two constants}
\label{sec:constants}

The geometric criterion uses exactly two solution-class constants,
both allowed to depend on the data \((u_0,\nu)\) and on nothing else
from the later evolution:
\begin{align}
\Omega&=\Omega(u_0,\nu)>0,
\label{eq:Omega}\\
C&=C(u_0,\nu)<\infty.
\label{eq:C}
\end{align}
Here \(\Omega\) is the vorticity-intensity threshold: \(\xi\) is constrained
only on \(\{r\ge\Omega\}\). The complementary set \(\{r<\Omega\}\) is
estimated without forming \(\xi\). The constant \(C\) is the uniform
spatial \(\tfrac12\)-H\"older modulus of \(\xi\) on that high set.
From~\eqref{eq:Omega}--\eqref{eq:C} the enstrophy argument produces one
derived majorant coefficient, depending on the two constants and on
viscosity,
\begin{equation}
\label{eq:Kconst}
K(C,\Omega,\nu)=\frac{C_\ast}{\nu}\,(C+\Omega)^2,
\end{equation}
where \(C_\ast<\infty\) depends only on the Sobolev, Hardy--Littlewood--Sobolev,
and Calder\'on--Zygmund constants of \(\T^3\) (or \(\R^3\)). The pair
\((\Omega,C)\) is the input; \(K\) is not an independent hypothesis.

These two constants are not constructed from~\eqref{eq:NS} in this note.
Theorem~\ref{thm:A} takes~\eqref{eq:Omega}--\eqref{eq:C} as given.

\section{The criterion}

\begin{theorem}[Beir\~ao da Veiga--Berselli, recast]
\label{thm:A}
Let \(u\) be the classical solution of~\eqref{eq:NS}--\eqref{eq:div} on
\([0,T^*)\), and let \(\Omega=\Omega(u_0,\nu)\) and \(C=C(u_0,\nu)\) be
the two constants~\eqref{eq:Omega}--\eqref{eq:C}. Suppose that for all
\(t\in[0,T^*)\) and all \(x,y\in\T^3\) with \(r(x,t)\ge\Omega\) and
\(r(y,t)\ge\Omega\),
\begin{equation}
\label{eq:holder}
\abs{\xi(x,t)-\xi(y,t)}\le C\,\abs{x-y}^{1/2}.
\end{equation}
Then \(T^*=\infty\), and \(u\in C^\infty(\T^3\times[0,\infty))\).
The continuation bound depends on the two constants only through \(K(C,\Omega,\nu)\)
in~\eqref{eq:Kconst}:
\begin{equation}
\label{eq:bound-thm}
\sup_{t<T^*}\norm{\omega(t)}_{L^2}^2
\le \norm{\omega(0)}_{L^2}^2
\exp\Bigl(\frac{K(C,\Omega,\nu)}{2\nu}\,\norm{u_0}_{L^2}^2\Bigr).
\end{equation}
\end{theorem}

The H\"older condition may be localized to \(\abs{x-y}<\delta\) for an
arbitrary \(\delta>0\); the far field is controlled by the integrable
periodized kernel. Only pairs in which \emph{both} magnitudes meet the
threshold \(\Omega\) are constrained. Pairs with \(r(x)<\Omega\) or
\(r(y)<\Omega\) are estimated crudely, without using \(\xi\).

\begin{remark}
On \(\{r(x),r(y)>0\}\) one has
\(\abs{\xi(x)-\xi(y)}=2\abs{\sin(\theta/2)}\) and
\(\sin\theta\le\abs{\xi(x)-\xi(y)}\), where \(\theta=\angle(\omega(x),\omega(y))\).
Thus~\eqref{eq:holder} implies
\(\sin\theta(x,y,t)\le C\abs{x-y}^{1/2}\) on the high-high set, which is
Assumption~A of Beir\~ao da Veiga--Berselli~\cite{BdVB2002} at exponent
\(\alpha=1/2\) with \(g\) a constant. Lipschitz direction
(Constantin--Fefferman~\cite{CF1993}) is the case \(\alpha=1\).
\end{remark}

\section{Enstrophy and stretching}

Take the \(L^2\) inner product of~\eqref{eq:vort} with \(\omega\). After
an integration by parts (justified on every compact subinterval of
\([0,T^*)\) by classical smoothness),
\begin{equation}
\label{eq:enstrophy}
\frac12\frac{\dd}{\dd t}\norm{\omega}_{L^2}^2+\nu\norm{\nabla\omega}_{L^2}^2
=\int_{\T^3}\mathcal K(x,t)\,\dd x,
\end{equation}
where the stretching density is
\begin{equation}
\label{eq:K}
\mathcal K(x,t):=\bigl((\omega\cdot\nabla)u\cdot\omega\bigr)(x,t).
\end{equation}
On \(\{r=0\}\) one has \(\mathcal K=0\). On \(\{r>0\}\) one has
\(\mathcal K=r^2(S\xi\cdot\xi)\).

Biot--Savart on \(\T^3\) reconstructs \(u\) from \(\omega\) by the
periodized Newtonian potential \(G\). Differentiating under the principal
value gives \(S\) as a Calder\'on--Zygmund operator of \(\omega\). The
Constantin--Fefferman representation~\cite{CF1993} then yields the pointwise
bound, for a dimensional constant \(c_0\) independent of \(u\),
\begin{equation}
\label{eq:Krep}
\abs{\mathcal K(x)}\le c_0\int_{\T^3}\frac{r(x)^2\,r(y)\,\sin\theta(x,y)}{\abs{x-y}^3}\,\dd y,
\end{equation}
where the integrand is read as \(0\) if \(r(x)=0\) or \(r(y)=0\), so that
\(\xi\) is never invoked on the degeneracy set. (On the torus, \(\abs{x-y}\)
denotes chordal/periodized distance; the local singularity is the same as
in \(\R^3\).)

\section{High-high versus complementary pairs}

Fix \(\Omega>0\) as in Theorem~\ref{thm:A}. Split the integral
in~\eqref{eq:Krep} according to whether \(r(y)\ge\Omega\) or not.
Write \(\mathcal K=\mathcal K_{\mathrm{hh}}+\mathcal K_{\mathrm{c}}\).

\paragraph{High-high.}
If \(r(x)\ge\Omega\) and \(r(y)\ge\Omega\),~\eqref{eq:holder} and
\(\sin\theta\le\abs{\xi(x)-\xi(y)}\) give
\(\sin\theta/\abs{x-y}^3\le C\abs{x-y}^{-5/2}\). Hence
\begin{equation}
\label{eq:Khh}
\abs{\mathcal K_{\mathrm{hh}}(x)}\le c_0 C\,r(x)^2\,I(x),
\qquad
I(x):=\int_{\T^3}\frac{r(y)}{\abs{x-y}^{5/2}}\,\dd y.
\end{equation}
The field \(I\) is a Riesz potential of order \(\beta=1/2\) applied to \(r\).
Hardy--Littlewood--Sobolev on \(\T^3\) (or \(\R^3\)) yields
\begin{equation}
\label{eq:HLS}
\norm{I}_{L^3}\le C_{\mathrm{HLS}}\norm{\omega}_{L^2},
\end{equation}
because \(1/3=1/2-(1/2)/3\). If \(r(x)<\Omega\), the factor \(r(x)^2\)
already places \(\mathcal K_{\mathrm{hh}}(x)\) in the complementary class
treated next; alternatively one may restrict~\eqref{eq:Khh} to
\(\{r(x)\ge\Omega\}\) and estimate the remainder with \(\sin\theta\le 1\)
and \(r(x)<\Omega\).

\paragraph{Complementary pairs.}
If \(r(y)<\Omega\), use \(\sin\theta\le 1\) and do \emph{not} form \(\xi(y)\).
After the usual cancellation in the CZ kernel of \(S\) (mean-zero
homogeneous kernel of degree \(-3\)), the complementary stretching is
controlled by the Hardy--Littlewood maximal function of the truncated
vorticity:
\begin{equation}
\label{eq:Kc}
\abs{\mathcal K_{\mathrm{c}}(x)}\le c_1 r(x)^2\,M\bigl(r\,1_{\{r<\Omega\}}\bigr)(x).
\end{equation}
In particular, using \(r1_{\{r<\Omega\}}\le\Omega\) and the boundedness of
\(M:L^2\to L^2\),
\begin{equation}
\label{eq:Kcint}
\int_{\T^3}\abs{\mathcal K_{\mathrm{c}}}\,\dd x
\le c_2\Omega\,\norm{\omega}_{L^2}\norm{\omega}_{L^6}.
\end{equation}
(The same bound follows from \(\abs{S(1_{\{r<\Omega\}}\omega)}\lesssim
M(r1_{\{r<\Omega\}})\). No division by \(r\) occurs.)

\section{Closing the enstrophy inequality}

Combining~\eqref{eq:enstrophy},~\eqref{eq:Khh},~\eqref{eq:HLS} and~\eqref{eq:Kcint},
\[
\int\abs{\mathcal K_{\mathrm{hh}}}\,\dd x
\le c_0 C\int r^2 I\,\dd x
\le c_0 C\norm{\omega}_{L^3}^2\norm{I}_{L^3}
\le c_3 C\norm{\omega}_{L^2}\norm{\omega}_{L^6}\norm{\omega}_{L^2},
\]
where the interpolation \(\norm{\omega}_{L^3}^2\le\norm{\omega}_{L^2}\norm{\omega}_{L^6}\)
was used. Thus
\begin{equation}
\label{eq:enst2}
\frac12\frac{\dd}{\dd t}\norm{\omega}_{L^2}^2+\nu\norm{\nabla\omega}_{L^2}^2
\le c_4(C+\Omega)\,\norm{\omega}_{L^2}^2\norm{\omega}_{L^6}.
\end{equation}
Sobolev on \(\T^3\) gives \(\norm{\omega}_{L^6}\le C_{\mathrm{S}}\norm{\nabla\omega}_{L^2}\).
Young's inequality, for every \(\varepsilon>0\),
\[
\norm{\omega}_{L^2}^2\norm{\omega}_{L^6}
\le \frac{\varepsilon}{c_4(C+\Omega)}\norm{\omega}_{L^6}^2
+\frac{c_4(C+\Omega)}{4\varepsilon}\norm{\omega}_{L^2}^4,
\]
so choosing \(\varepsilon\) small enough to absorb the first term into
\(\nu\norm{\nabla\omega}_{L^2}^2\) produces
\begin{equation}
\label{eq:gronwall}
\frac{\dd}{\dd t}\norm{\omega}_{L^2}^2
+\nu\norm{\nabla\omega}_{L^2}^2
\le K(C,\Omega,\nu)\,\norm{\omega}_{L^2}^4,
\end{equation}
with \(K(C,\Omega,\nu)\) the derived coefficient~\eqref{eq:Kconst}. The
only data-dependent constants entering \(K\) are the two constants
\(\Omega\) and \(C\).

\section{Continuation}

Energy equality for~\eqref{eq:NS} on \([0,T^*)\) yields
\begin{equation}
\label{eq:energy}
\frac12\norm{u(t)}_{L^2}^2+\nu\int_0^t\norm{\nabla u(s)}_{L^2}^2\,\dd s
=\frac12\norm{u_0}_{L^2}^2,
\end{equation}
hence \(\int_0^{t}\norm{\omega(s)}_{L^2}^2\,\dd s\le (2\nu)^{-1}\norm{u_0}_{L^2}^2\)
for every \(t<T^*\). Drop the dissipation in~\eqref{eq:gronwall} and apply
Gronwall:
\begin{equation}
\label{eq:bound}
\norm{\omega(t)}_{L^2}^2
\le \norm{\omega(0)}_{L^2}^2
\exp\Bigl(K\int_0^t\norm{\omega(s)}_{L^2}^2\,\dd s\Bigr)
\le \norm{\omega(0)}_{L^2}^2
\exp\bigl(K(2\nu)^{-1}\norm{u_0}_{L^2}^2\bigr).
\end{equation}
The right-hand side is independent of \(t<T^*\). Therefore
\(\sup_{t<T^*}\norm{\nabla u(t)}_{L^2}<\infty\). The blow-up
alternative~\eqref{eq:H1blowup} forces \(T^*=\infty\).

Once the solution exists globally in \(H^1\) with uniformly bounded
enstrophy, standard parabolic bootstrapping for~\eqref{eq:NS} (Kato,
or \(L^p\)-\(L^q\) estimates for the Stokes semigroup) upgrades
\(u\) to \(C^\infty(\T^3\times[0,\infty))\).

This is Theorem~\ref{thm:A}.

\section{Construction from the PDE, and the test}
\label{sec:test}

The two constants are now built from~\eqref{eq:NS}--\eqref{eq:div} and
from \(u_0\), then tested against the energy identity. No term is added
to the PDE. Degeneracy of \(\xi\) on \(\{r=0\}\) is respected.

\subsection{Irrotational data}

If \(\omega_0\equiv 0\), then \(u_0\) is a harmonic divergence-free field
on \(\T^3\), hence constant. The solution of~\eqref{eq:NS}--\eqref{eq:div}
is \(u(\,\cdot\,,t)=u_0\), which is smooth for all time. Take
\(\Omega=1\) and \(C=0\); the high set is empty. This case is settled
and is excluded below.

\subsection{The threshold \(\Omega\)}

Fix the intensity threshold independently of the later evolution,
\begin{equation}
\label{eq:Omega-construct}
\Omega(u_0,\nu):=1.
\end{equation}
Any other fixed positive number may replace \(1\); Beir\~ao da Veiga--Berselli
allow an arbitrary \(k>0\). The choice must \emph{not} track
\(\norm{\omega(t)}_\infty\). In particular one cannot take
\(\Omega(t)>\norm{\omega(t)}_\infty\), which would empty the high set
and make~\eqref{eq:holder} vacuous by circular use of a bound on
\(\omega\). The value~\eqref{eq:Omega-construct} depends on \((u_0,\nu)\)
only in that \(\nu>0\) and \(u_0\) is given; it is a universal positive
threshold.

\subsection{The modulus \(C(t)\) along the classical solution}

Let \(E(t):=\{x\in\T^3:r(x,t)\ge\Omega\}\). On each compact subinterval
of \([0,T^*)\) the solution is smooth, so \(E(t)\) is closed and \(\xi(t)\)
is smooth on a neighbourhood of \(E(t)\) whenever \(E(t)\) is nonempty
(because \(r\ge 1\) there). Define the running \(\tfrac12\)-H\"older
modulus on the high set by
\begin{equation}
\label{eq:C-running}
C(t):=
\begin{cases}
0 & \text{if }E(t)=\emptyset,\\
\sup\bigl\{\frac{\abs{\xi(x,t)-\xi(y,t)}}{\abs{x-y}^{1/2}}:
x,y\in E(t),\; x\neq y\bigr\}
& \text{if }E(t)\neq\emptyset.
\end{cases}
\end{equation}
Local smoothness gives \(C(t)<\infty\) for every \(t<T^*\). The constant
required by Theorem~\ref{thm:A} is the \emph{uniform} value
\begin{equation}
\label{eq:C-construct}
C(u_0,\nu):=\sup_{t\in[0,T^*)}C(t)\in[0,\infty].
\end{equation}
This is constructed from the PDE in the only non-circular sense available:
it is the actual modulus of the classical solution on its maximal interval.
Theorem~\ref{thm:A} applies if and only if this supremum is finite.

\subsection{Direction equation on \(\{r>0\}\)}

On \(\{r>0\}\), \(\abs{\xi}=1\) and \((\omega\cdot\nabla)u=S\omega\).
The chain rule in~\eqref{eq:vort} yields
\begin{equation}
\label{eq:xi-evol}
D_t\xi
=P_\xi(S\xi)
+\nu\, r^{-1}P_\xi(\Delta\omega),
\end{equation}
where \(D_t=\partial_t+u\cdot\nabla\) and
\(P_\xi v:=v-(\xi\cdot v)\xi\). Expanding \(\Delta\omega=\Delta(r\xi)\)
gives the equivalent viscous form
\begin{equation}
\label{eq:xi-visc}
r^{-1}P_\xi(\Delta\omega)
=\Delta\xi+2(\nabla\log r)\cdot\nabla\xi,
\end{equation}
in which \(\xi\cdot\Delta\xi=-\abs{\nabla\xi}^2\). The logarithmic
gradient is used only on \(\{r\ge\Omega\}\), where
\(\abs{\nabla\log r}\le\Omega^{-1}\abs{\nabla\omega}\). Thus on \(E(t)\)
\begin{equation}
\label{eq:gradxi}
\abs{\nabla\xi}\le\Omega^{-1}\abs{\nabla\omega}.
\end{equation}
No division by \(r\) occurs on \(\{r=0\}\).

\subsection{Morrey reduction}

Morrey's embedding on \(\T^3\) gives \(W^{1,6}\hookrightarrow C^{0,1/2}\).
Let \(\chi\in C^\infty(\R)\) satisfy \(\chi(s)=0\) for \(s\le 1\) and
\(\chi(s)=1\) for \(s\ge 2\), and set
\(v(x,t):=\chi\bigl(r(x,t)/\Omega\bigr)\,\xi(x,t)\) on \(\{r>0\}\),
with \(v=0\) on \(\{r=0\}\). Then \(v\) is a globally defined smooth
field on \([0,T^*)\), \(v=\xi\) on \(\{r\ge 2\Omega\}\), and
\begin{equation}
\label{eq:gradv}
\abs{\nabla v}\le C_\chi\,\Omega^{-1}\abs{\nabla\omega}.
\end{equation}
Hence a finite bound
\begin{equation}
\label{eq:morrey}
\sup_{t<T^*}\norm{\nabla\omega(t)}_{L^6(\T^3)}<\infty
\end{equation}
would imply \(\sup_{t<T^*}C(t)<\infty\) after replacing \(\Omega\) by
\(2\Omega\) in Theorem~\ref{thm:A}. Lipschitz control of \(\xi\) would
require \(\nabla\omega\in L^\infty_t L^\infty_x\), which is Beale--Kato--Majda
and is strictly stronger than~\eqref{eq:morrey}.

\subsection{What the energy identity supplies}

Energy equality~\eqref{eq:energy} yields only
\[
u\in L^\infty(0,T^*;L^2)\cap L^2(0,T^*;H^1),
\qquad
\omega\in L^2(0,T^*;L^2).
\]
It does \emph{not} yield \(\nabla\omega\in L^\infty_t L^6\), nor even
\(\nabla\omega\in L^2_t L^2\). The latter would be an enstrophy bound,
which is what Theorem~\ref{thm:A} is designed to produce after \(C\) is
known to be finite. Using~\eqref{eq:morrey} as a hypothesis to bound \(C\)
is therefore circular at the energy level.

A differential inequality for \(C(t)\) from~\eqref{eq:xi-evol} contains
the stretching \(P_\xi(S\xi)\) and the viscous term
\(\nu\Omega^{-1}P_\xi(\Delta\omega)\) on \(E(t)\). The strain \(S\) is a
Calder\'on--Zygmund image of \(\omega\). Neither term is controlled in
\(L^\infty_t L^6\) by \(\norm{\omega}_{L^2_t L^2_x}\) alone. The same
obstruction persists if \(\Omega=1\) is replaced by any other fixed
positive threshold depending only on \(\norm{u_0}_{H^s}\) and \(\nu\).

\subsection{Result of the test}

\begin{theorem}[construction from the PDE does not close]
\label{thm:B}
Let \(u_0\in C^\infty(\T^3;\R^3)\) be divergence-free and \(\nu>0\), and
let \(u\) be the classical solution of~\eqref{eq:NS}--\eqref{eq:div} on
\([0,T^*)\). Define \(\Omega\) by~\eqref{eq:Omega-construct} and
\(C(u_0,\nu)\) by~\eqref{eq:C-construct}. Then:
\begin{enumerate}
\item If \(\omega_0\equiv 0\), then \(T^*=\infty\) and \(u\) is globally
smooth.
\item If \(\omega_0\not\equiv 0\), then \(C(t)<\infty\) for every
\(t<T^*\), but the energy identity for~\eqref{eq:NS} does not imply
\(C(u_0,\nu)<\infty\). In particular it does not imply~\eqref{eq:morrey}.
\item Consequently Theorem~\ref{thm:A} cannot be applied from the energy
class alone. The test does not prove \(T^*=\infty\), and it does not
produce a finite-time singularity.
\end{enumerate}
\end{theorem}

Local smoothness on compact subintervals of \([0,T^*)\) always gives a
finite modulus on each such interval. Uniformity up to \(T^*\) is the
whole problem. Values \(\beta<1/2\) are not known to suffice
\cite{BdV2016}.

No term was added to~\eqref{eq:NS}. Degeneracy of \(\xi\) on \(\{r=0\}\)
was respected throughout. The Clay problem, as stated by
Fefferman~\cite{Fefferman2000}, remains open.

\begin{thebibliography}{9}
\bibitem{BdVB2002}
H.\ Beir\~ao da Veiga and L.\,C.\ Berselli,
\emph{On the regularizing effect of the vorticity direction in incompressible viscous flows},
Differential Integral Equations \textbf{15} (2002), 345--356.

\bibitem{BdV2016}
H.\ Beir\~ao da Veiga,
\emph{Open problems concerning the H\"older continuity of the direction of vorticity for the Navier--Stokes equations},
arXiv:1604.08083.

\bibitem{CF1993}
P.\ Constantin and C.\ Fefferman,
\emph{Direction of vorticity and the problem of global regularity for the Navier--Stokes equations},
Indiana Univ.\ Math.\ J.\ \textbf{42} (1993), 775--789.

\bibitem{Fefferman2000}
C.\ Fefferman,
\emph{Existence and smoothness of the Navier--Stokes equation},
Clay Mathematics Institute Millennium Problems, 2000.
\end{thebibliography}

\noindent\textit{License.} CC-BY-4.0.

\end{document}
